\documentclass[11pt,a4paper]{article}
\usepackage[utf8]{inputenc}
\usepackage[T1]{fontenc}
\usepackage{amsmath,amsfonts,amssymb}
\usepackage{graphicx}
\usepackage{hyperref}
\usepackage{listings}
\usepackage{color}
\usepackage{booktabs}
\usepackage{float}
\usepackage{geometry}
\geometry{margin=1in}
\definecolor{zenblue}{RGB}{41,121,255}
\definecolor{zengreen}{RGB}{52,199,89}
\hypersetup{colorlinks=true,linkcolor=zenblue,urlcolor=zenblue,citecolor=zenblue}

\title{\textbf{Zen-Max: Maximum Capability via Mixture-of-Experts Architecture}\\
\large Technical Report v2025.03}
\author{Zen LM Research Team\\
\texttt{research@zenlm.org}}
\date{March 2025}

\begin{document}
\maketitle

\begin{abstract}
We present \textbf{Zen-Max}, the highest-capability model in the Zen family, implementing the
Zen MoDE (Mixture of Distilled Experts) architecture at maximum scale: 480 billion total
parameters activating 48 billion per forward pass through a sparse mixture-of-experts (MoE)
routing mechanism. Trained on 5.5 trillion tokens with a 200K token context window, Zen-Max
establishes new performance records across the Zen family on reasoning, scientific understanding,
mathematics, and code generation benchmarks: MMLU 92.7\%, MATH 89.3\%, HumanEval 94.8\%,
GPQA 78.2\%, and MMMU 87.6\%. The MoE design achieves dense-model reasoning quality at
substantially lower per-token inference cost, making frontier capability accessible for research
and enterprise deployment at manageable compute budgets.
\end{abstract}

\tableofcontents
\newpage

%% ─────────────────────────────────────────────────────────────────────────────
\section{Introduction}

Scaling language model capability to the frontier has historically required proportional
increases in inference compute, creating tension between maximizing quality and deployment
cost. Mixture-of-experts (MoE) architectures \cite{shazeer2017outrageously, fedus2022switch}
address this tension by decoupling total parameter count from per-token active parameters:
a large expert pool provides rich representational capacity while sparse routing activates
only a relevant subset for each token.

Zen-Max realizes this principle at our largest scale to date: 480B total parameters, 48B
active per token (10\% activation rate), with 128 experts per MoE layer and top-8 routing.
This design delivers performance comparable to much larger dense models while constraining
inference FLOPs to approximately 48B-equivalent. We make the following contributions:

\begin{itemize}
  \item The Zen MoDE architecture at 480B/48B scale, with 128 experts, fine-grained expert
        segmentation, and auxiliary-loss-free load balancing via a bias-adjustment mechanism.
  \item A 200K token context window enabled by extended RoPE and a chunked attention
        implementation optimized for memory efficiency.
  \item RLVR post-training on an extended suite of verifiable domains including formal
        mathematics, code, logic puzzles, scientific question answering, and multi-step
        tool-use trajectories.
  \item State-of-the-art benchmark performance: MMLU 92.7\%, MATH 89.3\%, HumanEval 94.8\%,
        GPQA Diamond 78.2\%, MMMU 87.6\%.
\end{itemize}

%% ─────────────────────────────────────────────────────────────────────────────
\section{Architecture}

\subsection{Overall Structure}

Zen-Max is a decoder-only transformer with alternating dense attention layers and MoE
feed-forward layers. The 96-layer network has MoE FFN in 80 of the 96 layers, with 16 dense
layers (every 6th layer) to maintain stable global representations across the routing sparsity.

\begin{table}[H]
\centering
\caption{Zen-Max architecture hyperparameters.}
\label{tab:arch}
\begin{tabular}{lc}
\toprule
\textbf{Hyperparameter} & \textbf{Value} \\
\midrule
Parameters (total)               & 480B \\
Parameters (active per token)    & 48B \\
Layers (total)                   & 96 \\
Dense layers                     & 16 (every 6th) \\
MoE layers                       & 80 \\
Experts per MoE layer            & 128 \\
Experts activated per token      & 8 (top-8 routing) \\
Shared expert per MoE layer      & 2 (always active) \\
Attention heads                  & 128 \\
KV heads (GQA)                   & 8 \\
Hidden dimension                 & 7{,}168 \\
Expert FFN dimension             & 2{,}048 (per expert) \\
Vocabulary size                  & 151{,}936 \\
Context length (training)        & 200{,}000 \\
Position encoding                & RoPE ($\theta = 10{,}000{,}000$) \\
Activation                       & SiLU \\
Normalization                    & RMSNorm \\
\bottomrule
\end{tabular}
\end{table}

\subsection{Fine-Grained Expert Architecture}

Zen-Max uses a fine-grained expert decomposition where each MoE layer contains 128 small
experts rather than fewer large experts. Each expert $e_i$ applies a SwiGLU FFN with
intermediate dimension 2,048, smaller than the 7,168 hidden dimension:

\begin{equation}
  \text{Expert}_i(x) = \left(\text{SiLU}(xW_{\text{gate},i}) \odot xW_{\text{up},i}\right) W_{\text{down},i}
\end{equation}

The top-8 routing selects the 8 highest-scoring experts per token. Additionally, 2 shared
experts per layer are always activated and added to the routed experts' outputs. This
``shared expert'' pattern ensures that universal representations (stopwords, punctuation,
formatting) are handled without burdening the routing mechanism.

\subsection{Auxiliary-Loss-Free Load Balancing}

Expert load imbalance is a critical failure mode in MoE training, causing some experts to
receive disproportionate tokens (routing collapse) while others atrophy. Rather than adding
an auxiliary load balancing loss that trades task performance for balance, Zen-Max employs
a bias-adjustment mechanism \cite{deepseekmoe2024}:

\begin{equation}
  s_{i,j} = \text{softmax}(x_i W_g)_j + b_j
\end{equation}

where $b_j$ is a per-expert bias updated after each training step based on the deviation of
expert load from the target. Experts that are overloaded receive a negative bias; underloaded
experts receive a positive bias. This maintains load balance without introducing conflicting
gradient signals.

\subsection{200K Context Window}

A 200K token context window requires careful management of attention memory. Zen-Max combines:

\begin{enumerate}
  \item Extended RoPE with $\theta = 10{,}000{,}000$ providing position distinguishability
        up to 200K tokens.
  \item Chunked attention with chunk size $c = 4{,}096$ tokens, computing attention in chunks
        to avoid materializing the full $200K \times 200K$ attention matrix.
  \item A 3-stage training curriculum: 8K $\to$ 64K $\to$ 200K, each stage consuming
        20B, 10B, and 5B tokens respectively.
\end{enumerate}

\subsection{Multi-Head Latent Attention}

For the attention sublayer, Zen-Max uses Multi-Head Latent Attention (MLA) \cite{liu2024deepseekv2}
which compresses the KV cache via low-rank projections:

\begin{align}
  c^{KV} &= x W^{DKV} \in \mathbb{R}^{d_c} \\
  K &= c^{KV} W^{UK} \in \mathbb{R}^{n_h \times d_h} \\
  V &= c^{KV} W^{UV} \in \mathbb{R}^{n_h \times d_h}
\end{align}

where $d_c = 512 \ll n_h \cdot d_h$. Only $c^{KV}$ is cached (512 floats per layer per
token), reducing the KV cache by approximately 26$\times$ relative to full MHA at equivalent
head count. For a 200K-token sequence at 96 layers in BF16: $96 \times 200{,}000 \times 512
\times 2 \approx 19.7$ GB, compared to $\sim$510 GB for full MHA.

%% ─────────────────────────────────────────────────────────────────────────────
\section{Training Methodology}

\subsection{Pretraining at Scale}

Zen-Max pretrains on 5.5T tokens across the same domain distribution as Zen-Pro (see
Table~\ref{tab:data}), with an increased allocation to mathematical and scientific content.

\begin{table}[H]
\centering
\caption{Zen-Max pretraining data (5.5T tokens total).}
\label{tab:data}
\begin{tabular}{lcc}
\toprule
\textbf{Domain} & \textbf{Tokens (B)} & \textbf{Fraction} \\
\midrule
Web text (quality-filtered)   & 2{,}035 & 37.0\% \\
Code (all languages)          &   825 & 15.0\% \\
Books and long-form           &   770 & 14.0\% \\
Scientific articles           &   605 & 11.0\% \\
Mathematics (formal+informal) &   550 & 10.0\% \\
Multilingual web              &   385 &  7.0\% \\
Curated reasoning chains      &   330 &  6.0\% \\
\midrule
Total                         & 5{,}500 & 100.0\% \\
\bottomrule
\end{tabular}
\end{table}

Training infrastructure: 2048 H100-SXM5 GPUs, expert parallelism $\times$16, tensor
parallelism $\times$8, pipeline parallelism $\times$16, data parallelism $\times$8. Total
training compute approximately $3.2 \times 10^{24}$ FLOPs.

\subsection{MoE-Specific Training Stability}

MoE models are prone to instability from routing discontinuities. Zen-Max applies:

\begin{itemize}
  \item Gradient clipping at norm 1.0 with monitoring of per-expert gradient magnitudes.
  \item Expert dropout: randomly zeroing 1 expert per token during training for robustness.
  \item Router z-loss \cite{zoph2022stmoe}: $\mathcal{L}_z = \frac{1}{B}\sum_{x} (\log \sum_j e^{x_j})^2$
        applied at coefficient $10^{-3}$ to prevent logit explosion in the router.
\end{itemize}

\subsection{Post-Training: Extended RLVR}

Zen-Max applies RLVR on an expanded set of verifiable tasks beyond Zen-Pro's three domains:

\begin{itemize}
  \item Mathematics (AMC/AIME/Olympiad, formal Lean proofs)
  \item Code (Python, Rust, Go, C++ on competitive programming problems)
  \item Logic (first-order logic satisfiability, constraint satisfaction)
  \item Science QA (verified against authoritative databases: PubMed, arXiv theorem checks)
  \item Multi-step tool use (API call sequences verified by environment simulation)
\end{itemize}

GRPO sampling uses $G=16$ rollouts per prompt (doubled from Zen-Pro) to improve advantage
estimation quality at the 480B parameter scale.

%% ─────────────────────────────────────────────────────────────────────────────
\section{Evaluation}

\subsection{Core Benchmarks}

\begin{table}[H]
\centering
\caption{Zen-Max benchmark results versus frontier models.}
\label{tab:benchmarks}
\begin{tabular}{lcccc}
\toprule
\textbf{Benchmark} & \textbf{Zen-Max (480B)} & \textbf{Zen-Pro (72B)} & \textbf{Comp.\ A (large)} & \textbf{Comp.\ B (MoE)} \\
\midrule
MMLU (5-shot)          & \textbf{92.7} & 89.2 & 91.8 & 90.4 \\
MMLU-Pro               & \textbf{78.3} & 72.4 & 77.1 & 75.6 \\
ARC-Challenge          & 78.4 & 72.8 & 79.1 & 77.3 \\
HellaSwag              & 91.2 & 88.4 & 90.8 & 89.6 \\
\midrule
MATH (4-shot, CoT)     & \textbf{89.3} & 84.1 & 87.9 & 86.4 \\
AIME 2024 (pass@1)     & \textbf{66.7} & 53.3 & 63.3 & 60.0 \\
GSM8K                  & \textbf{96.2} & 93.4 & 95.8 & 94.7 \\
GPQA Diamond           & \textbf{78.2} & 71.8 & 76.4 & 74.1 \\
\midrule
HumanEval (pass@1)     & \textbf{94.8} & 91.3 & 93.2 & 91.8 \\
MBPP (pass@1)          & \textbf{88.4} & 84.6 & 86.9 & 85.3 \\
SWE-bench Verified     & \textbf{52.1} & 38.7 & 49.3 & 46.8 \\
\midrule
MMMU (val)             & \textbf{87.6} & --   & 86.1 & 84.3 \\
MT-Bench               & \textbf{9.4}  &  9.1 &  9.3 &  9.2 \\
\bottomrule
\end{tabular}
\end{table}

\subsection{Mathematical Olympiad Performance}

\begin{table}[H]
\centering
\caption{Performance on mathematical competition benchmarks.}
\label{tab:math}
\begin{tabular}{lcc}
\toprule
\textbf{Benchmark} & \textbf{Zen-Max} & \textbf{Prior Best} \\
\midrule
AIME 2024 (pass@1)   & 66.7\% & 63.3\% \\
AIME 2024 (pass@32)  & 93.3\% & 90.0\% \\
AIME 2023            & 71.4\% & 68.6\% \\
AMC 10               & 94.2\% & 92.8\% \\
AMC 12               & 91.8\% & 89.4\% \\
MATH Level 5 only    & 83.1\% & 80.7\% \\
OlympiadBench        & 67.4\% & 64.2\% \\
\bottomrule
\end{tabular}
\end{table}

\subsection{Long-Context Performance}

\begin{table}[H]
\centering
\caption{RULER scores and needle-in-a-haystack (NIAH) accuracy at increasing context lengths.}
\label{tab:longctx}
\begin{tabular}{lcccccc}
\toprule
\textbf{Model} & \textbf{8K} & \textbf{32K} & \textbf{64K} & \textbf{128K} & \textbf{200K} \\
\midrule
Zen-Max (480B)     & 97.1 & 95.3 & 92.7 & 88.4 & 82.1 \\
Zen-Pro (72B)      & 96.2 & 92.1 & 88.3 & 82.7 & --- \\
Competitor A       & 96.8 & 93.4 & 87.1 & 71.3 & 54.8 \\
\bottomrule
\end{tabular}
\end{table}

\subsection{Inference Cost Analysis}

The MoE sparsity of Zen-Max provides substantial cost advantages over dense models at
comparable quality.

\begin{table}[H]
\centering
\caption{Comparative inference cost (A100-80GB cluster, BF16).}
\label{tab:inference_cost}
\begin{tabular}{lcccc}
\toprule
\textbf{Model} & \textbf{Active Params} & \textbf{GPUs (TP8)} & \textbf{tok/s (batch=1)} & \textbf{Relative cost} \\
\midrule
Zen-Max (480B MoE) & 48B  & 8$\times$ A100 & 28  & 1.0$\times$ \\
Dense-equivalent   & 72B  & 8$\times$ A100 & 18  & 1.6$\times$ \\
Dense-equivalent   & 120B & 16$\times$ A100 & 11  & 4.2$\times$ \\
\bottomrule
\end{tabular}
\end{table}

%% ─────────────────────────────────────────────────────────────────────────────
\section{Expert Specialization Analysis}

We analyze expert routing patterns to understand emergent specialization. After training,
we feed 1M tokens from each domain and measure the Kullback-Leibler divergence between
per-domain routing distributions and the uniform distribution:

\begin{equation}
  \text{Specialization}(e, d) = D_{\text{KL}}(P(e \mid x \in d) \| P(e))
\end{equation}

We find that:

\begin{itemize}
  \item Mathematical tokens activate a consistent set of 12--18 experts (out of 128) at
        rates 4$\times$ above average, suggesting strong specialization.
  \item Code tokens show stronger expert specialization by programming language than by
        task type (generation vs.\ debugging), with Python and C$+$$+$ activating distinct
        expert subsets.
  \item Natural language tokens show weaker specialization than code or math, consistent
        with the hypothesis that diverse language modeling requires broader expert coverage.
\end{itemize}

%% ─────────────────────────────────────────────────────────────────────────────
\section{Related Work}

Sparse mixture-of-experts for language models was introduced in Sparsely-Gated MoE
\cite{shazeer2017outrageously} and scaled in Switch Transformer \cite{fedus2022switch}
and GLaM \cite{du2022glam}. Load balancing has been addressed via auxiliary losses
\cite{fedus2022switch}, expert capacity buffers \cite{lepikhin2020gshard}, and more
recently bias-adjustment mechanisms \cite{deepseekmoe2024}. Multi-Head Latent Attention
for KV cache compression follows \cite{liu2024deepseekv2}.

Long-context MoE models face additional challenges in routing consistency across long
sequences. Our chunked attention and staged context extension follow best practices
established in recent long-context scaling work \cite{peng2023yarn, hsieh2024ruler}.

%% ─────────────────────────────────────────────────────────────────────────────
\section{Limitations}

Expert routing introduces non-determinism that complicates reproducibility. Load balancing
remains imperfect: a small fraction of experts (3--5\%) consistently receives below-average
routing probability despite bias adjustment. Inference requires communication-heavy expert
parallelism, increasing latency on multi-node deployments. The 200K context window, while
large, degrades for generation tasks beyond 64K output tokens. MoE training instability
necessitates careful monitoring and checkpoint selection.

%% ─────────────────────────────────────────────────────────────────────────────
\section{Conclusion}

Zen-Max demonstrates that the Zen MoDE architecture scales favorably to 480B/48B active
parameters, achieving the highest benchmark scores in the Zen family: MMLU 92.7\%, MATH
89.3\%, HumanEval 94.8\%, GPQA Diamond 78.2\%, and MMMU 87.6\%. The sparse MoE design
delivers these results at the per-token FLOPs budget of a 48B dense model, providing a
compelling capability/cost trade-off for research and enterprise deployments. Zen-Max
serves as the primary model for demanding scientific, mathematical, and multi-step reasoning
workloads in the Zen family ecosystem.

%% ─────────────────────────────────────────────────────────────────────────────
\begin{thebibliography}{99}

\bibitem{shazeer2017outrageously}
N.~Shazeer et al., ``Outrageously Large Neural Networks: The Sparsely-Gated Mixture-of-Experts
Layer,'' \textit{ICLR}, 2017.

\bibitem{fedus2022switch}
W.~Fedus, B.~Zoph, and N.~Shazeer, ``Switch Transformers: Scaling to Trillion Parameter Models
with Simple and Efficient Sparsity,'' \textit{JMLR}, 2022.

\bibitem{du2022glam}
N.~Du et al., ``GLaM: Efficient Scaling of Language Models with Mixture-of-Experts,''
\textit{ICML}, 2022.

\bibitem{lepikhin2020gshard}
D.~Lepikhin et al., ``GShard: Scaling Giant Models with Conditional Computation and Automatic
Sharding,'' \textit{ICLR}, 2021.

\bibitem{deepseekmoe2024}
DeepSeek AI, ``DeepSeek-V2: A Strong, Economical, and Efficient Mixture-of-Experts Language
Model,'' \textit{arXiv:2405.04434}, 2024.

\bibitem{liu2024deepseekv2}
A.~Liu et al., ``DeepSeek-V2,'' \textit{arXiv:2405.04434}, 2024.

\bibitem{zoph2022stmoe}
B.~Zoph et al., ``ST-MoE: Designing Stable and Transferable Sparse Expert Models,''
\textit{arXiv:2202.08906}, 2022.

\bibitem{peng2023yarn}
B.~Peng et al., ``YaRN: Efficient Context Window Extension,'' \textit{arXiv:2309.00071}, 2023.

\bibitem{hsieh2024ruler}
C.-Y.~Hsieh et al., ``RULER: What's the Real Context Size of Your Long-Context Language Models?''
\textit{arXiv:2404.06654}, 2024.

\bibitem{shao2024deepseekmath}
Z.~Shao et al., ``DeepSeekMath,'' \textit{arXiv:2402.03300}, 2024.

\bibitem{wei2022emergent}
J.~Wei et al., ``Emergent Abilities of Large Language Models,'' \textit{TMLR}, 2022.

\bibitem{lightman2023lets}
H.~Lightman et al., ``Let's Verify Step by Step,'' \textit{arXiv:2305.20050}, 2023.

\end{thebibliography}

\end{document}
